\section{Introduction}
\label{sec:intro}

\subsection{Free-energy computation along a reaction coordinate}
A central task in computational statistical mechanics is to compute the free
energy associated with a low-dimensional \emph{reaction coordinate} of a
high-dimensional system in thermal equilibrium~\citep{LelievreRoussetStoltz2010}.
For most of the development we take a configuration to be a point
$q=(x,y)\in\R^2$\footnote{We write a configuration as $q=(x,y)$ in the
two-dimensional cases; the many-body case of \cref{sec:case_wca} uses
$q\in\R^{2n}$ with a scalar collective variable $z(q)$ in place of $x$
(\cref{rem:manybody}). The symbol $q_t(\cdot)$, always carrying a time subscript
and a scalar argument, denotes instead the Fisher--Rao \emph{target marginal} of
\cref{sec:fr}; the two are never used in the same role.} drawn from the
canonical (Boltzmann--Gibbs) density
\begin{equation}
  \pi(x,y) \;\propto\; e^{-\beta V(x,y)},
  \label{eq:canonical}
\end{equation}
where $V$ is the potential energy and $\beta=1/(k_{\mathrm B}T)$ is the inverse
temperature. Fixing the reaction coordinate to the first coordinate,
$\xib(x,y)=x$, the free energy $F(x)$ and the local mean force $F'(x)$ along
$\xib$ are the objects of interest. The difficulty is metastability: when $V$
has well-separated basins along $x$, unbiased Langevin dynamics crosses the
barrier rarely, so a direct estimate of $F$ converges intolerably slowly.

\subsection{Adaptive Biasing Force}
The Adaptive Biasing Force (ABF) method~\citep{DarvePohorille2001,Henin2004,
Comer2015} removes this metastability adaptively. It maintains a running
estimate $\Fphat_t(x)$ of the mean force and applies the antiderivative
$\Bt(x)=\int^x\Fphat_t$ as a bias that depends only on $x$. As
$\Bt\to F$, the effective dynamics along $x$ becomes diffusive and the
reaction coordinate is explored freely, while --- crucially --- the bias does
not change the conditional law of the orthogonal degrees of freedom given $x$.
ABF is mathematically well understood, with long-time convergence
guarantees~\citep{LelievreRoussetStoltz2008}, and is a natural baseline against
which to measure any proposed enhancement.

\subsection{Why add a Fisher--Rao correction?}
ABF accelerates exploration through the \emph{force}, but it does not directly
control \emph{where the particles are}. Early in a simulation the estimated bias
is poor, so the empirical reaction-coordinate marginal $\phat_t(x)$ can remain
concentrated in the initial basin, starving other values of $x$ of the samples
ABF needs to estimate $F'(x)$ there. A Fisher--Rao (FR) birth--death
correction~\citep{LuLuNolen2019} addresses exactly this: interpreting the
evolution of $\phat_t$ as a gradient flow of relative entropy in the
Fisher--Rao geometry~\citep{Amari2016}, one reallocates probability mass by
\emph{killing} particles in over-represented regions and \emph{cloning}
particles in under-represented regions, driving $\phat_t$ toward a chosen target
marginal $q_t(x)$ without integrating any additional dynamics. Coupling such a
correction to ABF is appealing because the two act on complementary objects
(force versus density), but it is not automatically beneficial: a birth--death
step resamples \emph{whole} configurations $(x,y)$, and if it distorts the
conditional law $Y\mid X=x$ it will corrupt precisely the quantity ABF estimates.

\subsection{Central question and contributions}
This report asks how FR changes ABF across five core controlled systems --- and, in a
final study, along the collective variable itself --- not
whether one slogan can describe all of them:
\begin{enumerate}
  \item \textbf{When does marginal reallocation help ABF, and when can it hurt?}
        (Q1)
  \item \textbf{What target and rate choices make FR useful in each regime?}
        (Q2--Q3)
  \item \textbf{What limitations and failure modes appear when the system becomes
        more realistic?} (Q4)
  \item \textbf{On a genuinely molecular reaction coordinate (a signed dihedral),
        does marginal FR help, is it superfluous, or can it produce a ``false
        improvement'' that flattens the biased marginal while corrupting a hidden
        conditional law?} (Q5)
\end{enumerate}
The systems are (i) a two-dimensional double-well \emph{metastability} model
with a quadrature reference (\cref{sec:case_meta}); (ii) a two-dimensional
\emph{entropic-bottleneck} model with a narrow transverse channel, for which both
the reference free energy \emph{and} the conditional law $Y\mid X=x$ are analytic
(\cref{sec:case_eb}); (iii) a many-body \emph{Weeks--Chandler--Andersen
(WCA) dimer} in a dense solvent with a high-accuracy thermodynamic-integration
reference (\cref{sec:case_wca}); and two united-atom alkanes whose slow coordinate
is a signed dihedral (\cref{sec:alkane_model}) --- (iv) \emph{butane}, an easy
control with an exact analytic torsional reference (\cref{sec:case_butane}), and
(v) \emph{pentane}, whose two coupled dihedrals make the second a hidden slow
coordinate when only the first is biased (\cref{sec:case_pentane}). Together they
span an easy baseline where target direction can matter, a cold analytic bottleneck
where density correction and free-energy accuracy can come apart, a many-body system
where improved coverage must be balanced against lineage diversity, and two
molecular torsions that test the mechanism where the collective variable is a
nonlinear function of a three-dimensional configuration.

Our contributions are: (i) a consistent ABF--FR notation that covers the analytic
two-dimensional models, the many-body WCA collective variable, and the periodic
molecular dihedral with its correct nonlinear generalized mean force
(\cref{sec:abf,sec:fr,sec:alkane_force}); (ii) matched-seed production studies
that quantify both the practical estimated-target method and diagnostic
uniform/oracle targets, including a pre-registered practical-equivalence test for the
easy butane control and configuration-resolved hidden-conditional diagnostics for
pentane (\cref{sec:case_meta,sec:case_eb,sec:case_wca,sec:case_butane,sec:case_pentane});
and (iii) a cross-case synthesis (\cref{sec:synthesis}) that explains why the
conclusions are intentionally different across examples. A recurring methodological point is the careful
interpretation of an \emph{oracle} target that uses the reference free energy: it
is a diagnostic control, never a deployable method and not a universal upper
bound. Its behaviour is itself part of the evidence: best in the easy
metastability setting, approximately tied with the practical target in WCA, and
worse than ABF in the cold entropic bottleneck.
